\section{Intelligence and Security Context}
\label{section:11-Intelligence-Security-Context}

% War is fundamentally a human activity, long understood to have an enduring \textit{nature} derived from its human-centricity and variable \textit{character} reflective of the context in which it is fought.  
% It has long been characterized by the \textit{danger} of armed conflict, the \textit{uncertainty} induced by the fog of war, the \textit{friction} generated by operating in a contested environment and pure \textit{chance}~\cite{Clausewitz2003}.


Many intelligence and security problems can be addressed computationally... 


\subsection{Use Cases for Intelligence and Security Problems}
\label{section:IS-Use-Cases}

%Overview of use-cases here. 

\subsubsection{Link Analysis}
\label{section:Link-Analysis}
% Reachability / Shortest Path
A link analysis problem requires an intelligence analyst to determine if two entities, events or other phenomena are related. 
Frequently a question of social contact, an analyst will build a list of known associates for each entity in question and expand those lists outwards until identifying a connection between the two. 
The problem can be multidimensional, friends list on a social media platform may give one subset of connections, the call data logs for the entities may give another and the locations that emerge in their pattern of life may form a third. 
Traditionally approached with pictures and red twine, more recently applications like the Analyst's Notebook\footnote{\url{https://i2group.com/i2-analysts-notebook}} are used to create digital equivalents.
Asking \textit{``Does this target associate with that target''} is then a problem of following the twine to see where (if at all) the two are connected. 

\subsubsection{Kill Chain Analysis}
\label{section:Kill-Chain-Analysis}
% Reachability / Shortest Path (Physical, Cyber)
A kill chain analysis is a popular methodology for post hoc incident analysis. 
The general idea is that for an incident to occur, an ordered set of events must precede it to the conditions for the incident. 
The canonical example draws from the counter-terrorism literature of the early 2000s, which asserts that for a terrorist attack to occur an attacker must select the target, then reconnoiter it, select specific targets, surveil the target, rehearse the attack, conduct the attack and then exploit their successes and escape~\cite{Tradoc2007}.
By understanding the chain of events, intelligence analysts seek to identify indicators and warnings from `left of bang' that may indicate an attack is imminent. 
To prevent an attacker from exploiting their success, you would block their actions on the objective. 
To prevent the attack, you catch them in rehearsals. 
To stop them from rehearsing, you should catch them surveilling the target, and so on.
Kill chains have also been popularised in the cyber domain by Lockheed Martin~\cite{Hutchins2011}.
It follows the same idea that the goal of a cyber threat intelligence analyst should be to detect adversary activity long before they can act on their objective. 
The challenge facing the analysts is how to identify when these events, or strings of events are occurring as part of that kill chain, rather than being indicative of random noise.

\subsubsection{Call Chaining Analysis}
\label{section:Call-Chaining}
% Graph Search

\subsubsection{Intelligence Preparation of the Battlespace (Population Segmentation)}
\label{section:IPB}
% Community Detection / Graph Segmentation

\subsubsection{Entity Resolution}
\label{section:Entity-Resolution}
% Entity Resolution (Time Constrained)

\subsubsection{Event and Entity Deduplication}
\label{section:Event-Entity-Deduplication}
% Entity Resolution

\subsubsection{Co-Traveler Analysis}
\label{section:Co-Traveler-Analysis}
% Contact Tracing, Drug Interdiction etc
% Algorithm??

\subsubsection{Financial Analysis}
\label{section:Financial-Analysis}
% Subgraph Matching

\subsubsection{Link Prediction}
\label{section:Link-Prediction}
% GNN Missing Link Prediction

\subsubsection{Intrusion Detection}
\label{section:Intrusion-Detection}

\subsubsection{Course of Action Prediction}
\label{section:Course-Of-Action-Prediction}

%
% Civy world - Recommend a friend on a platform (i.e. id who is most suited to be the bridge node -- MIT research?)

\subsubsection{Object Detection and Image Classification}
\label{section:Object-Detection-Image-Classification}
% Does Graph Accelerator improve classical NN speeds? 

